\section*{Chapitre 14 - Pourcentage}
\addcontentsline{toc}{section}{Chapitre 14 - Pourcentage}



\subsection*{I. Proportion et pourcentage}
\addcontentsline{toc}{subsection}{I. Proportion et pourcentage.}

\begin{Exemple}
Un savon pesant 100 g contient 38 g d'huile.\\
On dit qu'un savon contient 38 g d'huile pour 100 g de savon. On le note : $38 \%$.
\end{Exemple}

\begin{Définitions}[c] $ $
\begin{itemize}[label=\textcolor{Couleur6}{$\bullet$}]
\item Pour tout nombre $a$, le nombre $\dfrac{a}{100}$ est appelé \textbf{pourcentage}. \\
On peut aussi l'écrire $a$ \%.
\item Une \textbf{proportion} exprime le rapport entre une partie et un tout.
\end{itemize}
\end{Définitions}

\begin{Exemple}
Dans un panier de fruits, $8$ pommes sur $50$ fruits sont rouges. \\
Quelle est la proportion de pommes rouges ?

\ligne

\end{Exemple}



\subsection*{II. Calculer mentalement}
\addcontentsline{toc}{subsection}{II. Calculer mentalement}

\begin{center}
\renewcommand{\arraystretch}{1.2}
\begin{tabular}{|>{\columncolor{Couleur10!10}}l|c|c|c|c|c|}
  \hline
Pourcentage  & 10\,\% & 25\,\% & 50\,\% & 75\,\% & 100\,\% \\
  \hline
  Prendre\ldots         & le dixième      & le quart        & la moitié       & les trois quarts & le tout \\
  \hline
  revient à\ldots       & $\div 10$       & $\div 4$        & $\div 2$        & $\div 4$ puis $\times 3$ & $\times 1$ \\
  \hline
\end{tabular}
\end{center}


\begin{Exemple}
\vspace{-1.5cm}
\begin{multicols}{3}

\begin{align*}
&50\% \text{ de } 60\text{ €}\\
&=  \ldots\ldots \\
&= \ldots\ldots
\end{align*}

\begin{align*}
&10\% \text{ de } 42\\
&= \ldots\ldots \\
&= \ldots\ldots
\end{align*}

\begin{align*}
&75\% \text{ de } 50 \text{ L}\\
&= \ldots\ldots\\
&= \ldots\ldots \\
&= \ldots\ldots
\end{align*}

\end{multicols}
\vspace{-0.5cm}
\end{Exemple}

\subsection*{III. Appliquer un pourcentage}
\addcontentsline{toc}{subsection}{III. Appliquer un pourcentage}

\begin{Règle}[c]
Calculer $p$ \% d'un nombre revient à multiplier ce nombre par $\dfrac{p}{100}$.
\end{Règle}

\begin{Exemples}
\begin{enumerate}
\item Calculer 30 \% de 80. 

\ligne




\item Un magasin propose une réduction de 15 \% sur un article qui coûte 60 €.\\ Quel est le montant de la réduction ?

\ligne

\ligne


\end{enumerate}
\end{Exemples}